\subsection{完备空间的子空间}

命题 8 完备空间的每个闭子空间都是完备的。豪斯多夫一致空间（完备与否）的每个完备子空间都是闭的。

设 X 是完备空间，A 是 X 的闭子空间。若 F 是 A 上的 Cauchy 滤子，则它是 X 上的 Cauchy 滤子基（第 1 目命题 3），因而收敛于某点 \(x \in X\)；但由于 A 是闭的，我们有 \(x \in A\)，因此 F 在子空间 A 中收敛。

现在设 A 是豪斯多夫一致空间 X 的非闭子集，且 \(b \in \overline{\mathbf{A}} - \mathbf{A}\)。b 在 X 中的邻域滤子 B 在 A 上的迹 \(\mathfrak{B}_{\mathbf{A}}\) 是 A 上的 Cauchy 滤子；但它不能收敛于某点 \(c \in \mathbf{A}\)，否则 c 将是 B 的极限点（第 2 目命题 5 推论 3），而这是荒谬的，因为 \(b \neq c\) 且 X 是豪斯多夫的。

命题 9 设 X 是一致空间，A 是 X 的稠密子集，且 A 上每个 Cauchy 滤子基都在 X 中收敛。则 X 是完备的。

只需证明 X 上每个极小 Cauchy 滤子 \(\mathfrak{F}\) 都收敛。由于 A 稠密，且 \(\mathfrak{F}\) 的每个集合都有非空内部（第 2 目命题 5 推论 4），迹 \(\mathfrak{F}_{\mathrm{A}}\)（\(\mathfrak{F}\) 在 A 上的）是 A 上的 Cauchy 滤子，因而收敛于某点 \(x_{0} \in \mathrm{X}\)。由于 \(\mathfrak{F}\) 粗于由 \(\mathfrak{F}_{\mathrm{A}}\) 在 X 上生成的滤子，故 \(\mathfrak{F}\) 收敛于 \(x_{0}\)（第 2 目命题 5 推论 3）。

